\chapter{Proyecto y diseño del sistema}
\label{cap:proyecto}

\section{Problema y propósito}
Un robot manipulador fijo puede orientar y posicionar una herramienta dentro de un espacio limitado por su geometría. Una base móvil amplía la región de trabajo, aunque introduce incertidumbre de localización y modifica la estabilidad del conjunto. AutoFinalRobo2 aborda esa integración mediante una plataforma Mecanum y un brazo de seis ejes.

El problema consiste en convertir una consigna de desplazamiento o una pose de herramienta en acciones coordinadas, manteniendo una correspondencia clara entre geometría, unidades, sentidos de giro, mediciones y restricciones físicas. La integración debe permitir explicar por qué gira cada rueda, cómo se obtiene una posición estimada, cómo cambian las coordenadas del efector y qué condiciones limitan su operación.

El objetivo general es desarrollar una plataforma móvil con capacidad de manipulación y documentar los principios que gobiernan su diseño y funcionamiento. Los objetivos específicos son:
\begin{itemize}
\item formular la cinemática de la base y la relación entre cuentas Hall y movimiento;
\item desarrollar las transformaciones, la cinemática y el Jacobiano del brazo;
\item describir motores, transmisiones, sensores, alimentación e interfaces;
\item relacionar estabilidad, masa y cargas con las posturas de trabajo;
\item comparar consignas y mediciones para interpretar el desempeño;
\item identificar las condiciones necesarias para completar la integración física.
\end{itemize}

\section{Alcance y metodología}
El sistema contempla movilidad planar, manipulación, operación local mediante pantalla y encoder, control desde un navegador y herramientas de mantenimiento. El diseño mecánico del brazo es la configuración v17 con eslabones de PVC de 25 y 20 mm de diámetro exterior y pinza de dos mandíbulas. La plataforma se apoya únicamente sobre sus ruedas.

El desarrollo se organiza en definición de requisitos, diseño CAD, modelado matemático, implementación y caracterización. Las dimensiones nominales representan geometría de diseño; las masas calculadas dependen de densidades y de la representación de componentes. Las mediciones describen únicamente las condiciones en que fueron obtenidas. Esta distinción permite formular conclusiones proporcionadas a la evidencia disponible.

\begin{table}[!htbp]
\centering\small
\caption{Requisitos funcionales y forma de abordarlos.}
\begin{tabularx}{\linewidth}{@{}P{33mm}YY@{}}
\toprule
\textbf{Necesidad} & \textbf{Solución del proyecto} & \textbf{Condición de aceptación}\\\midrule
Movilidad planar & Cuatro ruedas Mecanum y mezcla de velocidades. & Comparar recorrido y giro con una referencia externa.\\
Posicionar herramienta & Manipulador de seis ejes y definición de TCP. & Contrastar poses y respetar límites e interferencias.\\
Sujetar objetos & Pinza de dos cremalleras accionadas por un piñón. & Caracterizar fuerza de contacto y retención.\\
Operar localmente & OLED monocroma y encoder con pulsador. & Navegación consistente y cambios confirmados.\\
Operar a distancia & Web, Wi-Fi, ESP y enlace serie al ATmega. & Consignas compatibles y estado de ejecución observable.\\
Mantener estabilidad & Restricción de postura y análisis de apoyo. & Centro de masa y momentos con margen suficiente.\\
\bottomrule
\end{tabularx}
\end{table}

\section{Arquitectura general}
El ATmega2560 concentra el control de los motores de la plataforma, adquisición de sensores e interacción con la OLED y el encoder. El ESP-12S, basado en ESP8266, aporta conectividad, servidor HTTP, WebSocket y almacenamiento LittleFS. La aplicación React y Vite se ejecuta en el navegador y presenta estado, control, visualización y herramientas de exploración.

La arquitectura se completa con las definiciones compartidas de comunicación y unidades, los modelos CAD y un cargador de mantenimiento para Windows. El software de ATmega está escrito en C con avr-libc, sin entorno Arduino ni DMA; el ESP utiliza Arduino ESP8266. La OLED es el único dispositivo I2C y su acceso se centraliza en el I2C Manager.

\begin{figure}[!htbp]
\centering
\begin{tikzpicture}[node distance=8mm and 7mm,box/.style={draw=azul!35,fill=gris,rounded corners=2mm,text width=39mm,minimum height=17mm,align=center,font=\small},arr/.style={-{Stealth[length=2mm]},thick,granate}]
\node[box] (web) {Operador y Web\\Control y visualización};
\node[box,right=of web] (esp) {ESP8266\\Wi-Fi, HTTP y WS};
\node[box,right=of esp] (avr) {ATmega2560\\Control y adquisición};
\node[box,below=of web] (tools) {CAD y herramientas\\Diseño y mantenimiento};
\node[box,below=of esp] (local) {OLED y encoder\\Interfaz local};
\node[box,below=of avr] (io) {Motores y sensores\\Plataforma física};
\draw[arr] (web)--node[above=8mm,font=\scriptsize]{WebSocket}(esp);
\draw[arr] (esp)--node[above=8mm,font=\scriptsize]{UART}(avr);
\draw[arr] (avr)--(io);
\draw[arr] (avr.south west)--(local.north east);
\end{tikzpicture}
\caption{Responsabilidades funcionales de la plataforma. La potencia y el bus STS del brazo se integran como un subsistema específico.}
\label{fig:arquitectura}
\end{figure}
\FloatBarrier

\section{Configuración mecánica}
La plataforma emplea cuatro ruedas Mecanum accionadas individualmente. El brazo incorpora seis articulaciones rotacionales: giro de base, hombro, codo, giro de muñeca, inclinación de muñeca y giro final de herramienta. La pinza añade apertura, sin aumentar los seis componentes de la pose rígida del TCP.

\figura[.69]{brazo.png}{Representación CAD del manipulador y sus interfaces principales.}{fig:brazo}

\begin{table}[!htbp]
\centering\small
\caption{Parámetros nominales de geometría y estructura.}
\begin{tabularx}{\linewidth}{@{}P{51mm}P{27mm}Y@{}}
\toprule\textbf{Magnitud} & \textbf{Valor} & \textbf{Interpretación}\\\midrule
Diámetro de rueda & 80 mm & Geometría nominal de la plataforma.\\
Suma de semidistancias & 170 mm & Parámetro de mezcla Mecanum.\\
Hombro--codo & 70 mm & Distancia geométrica entre ejes.\\
Codo--J4 & 60 mm & Tramo distal al codo.\\
J4--J5 & 55 mm & Separación hasta el centro de muñeca.\\
J5--J6 & 65 mm & Tramo distal sobre el eje de muñeca.\\
J6--TCP & 78,5 mm & Desplazamiento hasta el punto de herramienta.\\
Tubo del brazo & $\varnothing 25\times1{,}2$ mm & Diámetro interior nominal 22,6 mm.\\
Tubo del antebrazo & $\varnothing 20\times1$ mm & Diámetro interior nominal 18 mm.\\
Longitudes de corte & 42 y 28,5 mm & Piezas de tubo; no son distancias entre ejes.\\
\bottomrule
\end{tabularx}
\end{table}

La estructura se divide en montura y control fijo, base y hombro, brazo, antebrazo y giro J4, muñeca J5--J6 y herramienta. El accionamiento del codo está próximo al hombro y transmite por correa. Al acercar una masa al eje se reduce su contribución gravitatoria, aunque el cálculo completo también debe incluir poleas, soportes, tensión e inercias.

Se seleccionan cinco servos STS3250 para J1 a J5 y dos STS3032 A0 para J6 y pinza. Las reducciones incrementan el par de salida y disminuyen la velocidad. La forma CAD define la ocupación geométrica; ajustes, roscas, tolerancias, orientación de impresión y referencias comerciales deben cerrarse antes de liberar fabricación.

\subsection{Transmisiones y selección de materiales}
\begin{table}[!htbp]
\centering\small
\caption{Transmisiones nominales del manipulador.}
\begin{tabularx}{\linewidth}{@{}P{19mm}YP{28mm}@{}}
\toprule\textbf{Eje} & \textbf{Transmisión} & \textbf{Relación motor/salida}\\\midrule
J1 & Directa. & 1\\
J2 & Piñón de 18 y rueda de 72 dientes, módulo 0,8. & 4\\
J3 & Correa HP3M $210\times6$ y poleas de 20/50 dientes. & 2,5\\
J4 & Engranajes de 24/24 dientes. & 1\\
J5 y J6 & Directas. & 1\\\bottomrule
\end{tabularx}
\end{table}
Los tubos de PVC ofrecen secciones cerradas con masa reducida. Las uniones impresas resuelven alojamientos y enlaces entre piezas metálicas, rodamientos y actuadores. La resistencia de una pieza FDM depende del material efectivo, orientación, paredes, relleno, temperatura y fijación; una densidad de CAD no sustituye esas propiedades.

\subsection{Fabricación de las piezas por impresión 3D}
Fabricamos por impresión 3D FDM las piezas propias del brazo, como carcasas, alojamientos, soportes, uniones y elementos de transmisión. Los tubos de PVC, rodamientos, tornillería y actuadores se incorporan como componentes separados. El porcentaje de relleno interior, o \emph{infill}, se elige según la función mecánica de cada pieza: se reduce en las envolventes y aumenta donde deben transmitirse fuerzas y momentos.

\begin{table}[!htbp]
\centering\small
\caption{Criterio de relleno de las piezas impresas del brazo.}
\begin{tabularx}{\linewidth}{@{}P{37mm}P{20mm}Y@{}}
\toprule\textbf{Grupo de piezas} & \textbf{Relleno} & \textbf{Aplicación y motivo}\\\midrule
Carcasas y alojamientos no estructurales & 25\,\% & Todos los \emph{housings}, tapas, cubiertas y secciones no cargadas: dan forma, protegen y alojan componentes con menor consumo de material.\\
Estructura y soportes cargados & 80\,\% & Torre entre J1 y J2, soportes estructurales y uniones que transmiten momentos: se aumenta el relleno donde pasa la carga del brazo.\\
Engranajes y transmisión impresa & 100\,\% & Engranajes, piñones, cremalleras y demás piezas impresas de transmisión que requieren construcción maciza.\\\bottomrule
\end{tabularx}
\end{table}

\needspace{5\baselineskip}
La clasificación depende del recorrido de la carga. Una carcasa se imprime al 25\,\% cuando cumple una función de envolvente; si una zona también actúa como soporte y recibe el momento de una articulación, se trata como sección estructural al 80\,\%. Los engranajes se configuran macizos, con relleno nominal del 100\,\%, para evitar un interior de baja densidad bajo los dientes y en el cuerpo que transmite el par.

El momento aumenta con la fuerza y con su distancia perpendicular al eje, $M=F\,d_{\perp}$. Por eso la torre J1--J2 y sus soportes necesitan mayor relleno: reciben el efecto del peso y de las acciones del conjunto situado por encima. El relleno se combina con paredes, orientación de capas y fijaciones; por sí solo no determina una capacidad de carga. Los porcentajes indicados son el criterio de fabricación adoptado y no un resultado de resistencia medida.

Una definición de pieza identifica geometría y una ocurrencia identifica una instancia montada. Las cantidades y masas deben sumarse por ocurrencia. Un disco de paso de engranaje representa una interfaz cinemática, pero no equivale a un perfil de diente listo para fabricar.
\FloatBarrier
